\section{Articulación con las Misiones de la Política de Investigación e Innovación}

El presente proyecto se articula principalmente con la Misión de Bioeconomía y Territorio, al enfocarse en la transformación y valorización de un recurso biológico estratégico para el país como lo es el café. En particular, la propuesta contribuye al desarrollo de soluciones tecnológicas orientadas a mejorar los procesos agroindustriales, promoviendo la generación de valor agregado y la competitividad del sector cafetero colombiano en mercados internacionales.

De manera específica, el proyecto permite avanzar en la diversificación de la oferta exportable de café, al facilitar la estandarización y escalamiento de procesos de cafés diferenciados como honey y natural, los cuales presentan un alto potencial en mercados especializados. La implementación de tecnologías de secado basadas en el control de humedad relativa permite un manejo más preciso del proceso, favoreciendo la obtención de perfiles sensoriales diferenciados y consistentes, lo que contribuye a posicionar el café colombiano en segmentos de mayor valor agregado dentro de la bioeconomía.

Asimismo, la investigación se alinea con la Misión de Energía Sostenible, mediante la implementación de tecnologías de secado basadas en bomba de calor y la integración potencial de fuentes de energía renovable, como la energía solar. Este enfoque permite optimizar el consumo energético del proceso, reducir la dependencia de combustibles convencionales y avanzar hacia sistemas de producción más sostenibles.

Adicionalmente, el proyecto presenta una articulación con la Misión de Ciencia para la Paz, al contribuir al fortalecimiento de economías lícitas en territorios rurales a través del mejoramiento de la productividad y calidad del café. La implementación de tecnologías que faciliten el procesamiento centralizado del grano, en el marco de esquemas como el modelo CERPER, permite generar condiciones más equitativas para los caficultores, reducir riesgos asociados a la variabilidad en la calidad y promover ingresos más estables.

En este sentido, el proyecto aporta a la construcción de entornos productivos sostenibles, ofreciendo alternativas que fortalecen la permanencia en actividades agrícolas legales y contribuyen al desarrollo territorial, especialmente en regiones donde el café representa una opción estratégica frente a otras dinámicas productivas.

En conjunto, la propuesta se enmarca en los objetivos de la Política de Investigación e Innovación Orientada por Misiones, al integrar desarrollo tecnológico, sostenibilidad energética e impacto socioeconómico, contribuyendo a la modernización del sector cafetero colombiano y al desarrollo territorial.